% --- Archon physics formalization source begin ---
% archon:physics
% archon:covers PhyXMiniProblems/problem_phyx_mini_0819.lean
% archon:source-report reports/phyx_mini/problem_phyx_mini_0819.source.json

\chapter{Physics problem phyx\_mini\_0819}
\label{ch:PhyXMiniProblems_problem_phyx_mini_0819}

\paragraph{Problem source.}
\#\# Physical scenario

A machine part as shown in figure consists of three small disks linked by lightweight struts.

\#\# Question

What is the body's kinetic energy if it rotates about axis 1 with angular speed \textbackslash{}( \textbackslash{}omega = 4.0 \textbackslash{}text\{ rad/s\} \textbackslash{})?

\#\# Answer choices

- A: A: \textbackslash{}( 0.66 \textbackslash{}text\{ J\} \textbackslash{})

- B: B: \textbackslash{}( 0.46 \textbackslash{}text\{ J\} \textbackslash{})

- C: C: \textbackslash{}( 1.32 \textbackslash{}text\{ J\} \textbackslash{})

- D: D: \textbackslash{}( 0.92  \textbackslash{}text\{ J\} \textbackslash{})

\#\# Figure caption (auxiliary; use the image as primary evidence)

The image is a diagram showing three disks labeled A, B, and C. 

- **Disk A** is marked with an "Axis 1" perpendicular to the plane of the figure, has a mass \textbackslash{}( m\_A = 0.30 \textbackslash{}, \textbackslash{}text\{kg\} \textbackslash{}), and is connected to Disk B by a bar measuring \textbackslash{}( 0.50 \textbackslash{}, \textbackslash{}text\{m\} \textbackslash{}).
  
- **Disk B** has a mass \textbackslash{}( m\_B = 0.10 \textbackslash{}, \textbackslash{}text\{kg\} \textbackslash{}) and is connected directly below Disk C by a vertical bar measuring \textbackslash{}( 0.30 \textbackslash{}, \textbackslash{}text\{m\} \textbackslash{}). An "Axis 2" runs through disks B and C.
  
- **Disk C** has a mass \textbackslash{}( m\_C = 0.20 \textbackslash{}, \textbackslash{}text\{kg\} \textbackslash{}) and is connected to Disk A by a bar measuring \textbackslash{}( 0.40 \textbackslash{}, \textbackslash{}text\{m\} \textbackslash{}).

The lines connecting disks represent rigid bars or rods, forming a triangular configuration. The masses and distances are annotated next to the relevant sections of the diagram.

\#\# Dataset metadata

Rotational Motion; Physical Model Grounding Reasoning, Spatial Relation Reasoning

\paragraph{Recorded answer/context.}
B: B: \textbackslash{}( 0.46 \textbackslash{}text\{ J\} \textbackslash{})

\paragraph{Figure/image path.}
/root/proposal\_for\_physic/science-mango/phyx\_mini\_run/phyx\_data/test\_image/819.png

\paragraph{Formalization target.}
create a compiling Lean file with sorry bodies at `PhyXMiniProblems/problem\_phyx\_mini\_0819.lean`. The Lean declarations must preserve the physical quantities, dimensions or dimensional roles, figure labels, governing-law hypotheses, and final relation expressed by this problem.
Use Mathlib/Physlib names found through LeanExplore where available. If a physics API is missing, introduce faithful local abstractions rather than scalar placeholder aliases.

\begin{theorem}[Physics formalization target]
\label{thm:physics:phyx_mini_0819:target}
\lean{PhyXMiniProblems.ProblemPhyXMini0819.kineticEnergyAboutAxis1_is_answer_B}
\uses{def:physics:phyx-mini-0819:phyxminiproblems-problemphyxmini0819-energyinjoules, def:physics:phyx-mini-0819:phyxminiproblems-problemphyxmini0819-threediskmachinepart, def:physics:phyx-mini-0819:phyxminiproblems-problemphyxmini0819-matchesprimaryfigure, def:physics:phyx-mini-0819:phyxminiproblems-problemphyxmini0819-matchesproblemdata, def:physics:phyx-mini-0819:phyxminiproblems-problemphyxmini0819-satisfiesthreediskrotationallaws, def:physics:phyx-mini-0819:phyxminiproblems-problemphyxmini0819-displayedenergyjoules, def:physics:phyx-mini-0819:phyxminiproblems-problemphyxmini0819-recordedanswerchoice, def:physics:phyx-mini-0819:phyxminiproblems-problemphyxmini0819-roundstonearesthundredth, def:physics:phyx-mini-0819:phyxminiproblems-problemphyxmini0819-isuniqueclosestenergychoice}
The assigned autoformalize agent should translate this physics problem into Lean declarations in the covered file, with theorem and lemma proof bodies written as `by sorry`.
\end{theorem}
\begin{proof}
This is an autoformalization task, not a proof task. Produce statements that can later be proved by the physics prover without weakening the physical model.
\end{proof}

% --- Archon declaration topology begin ---
\section{Declaration topology}
\begin{definition}[Mass Quantity]
  \label{def:physics:phyx-mini-0819:phyxminiproblems-problemphyxmini0819-massquantity}
  \lean{PhyXMiniProblems.ProblemPhyXMini0819.MassQuantity}
  A nonnegative, unit-independent physical mass
\end{definition}
\begin{proof}
  This is the declared model definition of Mass Quantity.
\end{proof}
\begin{definition}[Length Quantity]
  \label{def:physics:phyx-mini-0819:phyxminiproblems-problemphyxmini0819-lengthquantity}
  \lean{PhyXMiniProblems.ProblemPhyXMini0819.LengthQuantity}
  A nonnegative, unit-independent physical length
\end{definition}
\begin{proof}
  This is the declared model definition of Length Quantity.
\end{proof}
\begin{definition}[Angular Speed Quantity]
  \label{def:physics:phyx-mini-0819:phyxminiproblems-problemphyxmini0819-angularspeedquantity}
  \lean{PhyXMiniProblems.ProblemPhyXMini0819.AngularSpeedQuantity}
  A nonnegative angular-speed magnitude; radians are dimensionless
\end{definition}
\begin{proof}
  This is the declared model definition of Angular Speed Quantity.
\end{proof}
\begin{definition}[Moment Of Inertia Quantity]
  \label{def:physics:phyx-mini-0819:phyxminiproblems-problemphyxmini0819-momentofinertiaquantity}
  \lean{PhyXMiniProblems.ProblemPhyXMini0819.MomentOfInertiaQuantity}
  A nonnegative scalar moment of inertia, with dimension M L²
\end{definition}
\begin{proof}
  This is the declared model definition of Moment Of Inertia Quantity.
\end{proof}
\begin{definition}[Energy Quantity]
  \label{def:physics:phyx-mini-0819:phyxminiproblems-problemphyxmini0819-energyquantity}
  \lean{PhyXMiniProblems.ProblemPhyXMini0819.EnergyQuantity}
  Rotational kinetic energy, with dimension M L² T⁻²
\end{definition}
\begin{proof}
  This is the declared model definition of Energy Quantity.
\end{proof}
\begin{definition}[mass In Kilograms]
  \label{def:physics:phyx-mini-0819:phyxminiproblems-problemphyxmini0819-massinkilograms}
  \lean{PhyXMiniProblems.ProblemPhyXMini0819.massInKilograms}
  \uses{def:physics:phyx-mini-0819:phyxminiproblems-problemphyxmini0819-massquantity}
  Read a physical mass in coherent-SI kilograms
\end{definition}
\begin{proof}
  This is the declared model definition of mass In Kilograms.
\end{proof}
\begin{definition}[length In Meters]
  \label{def:physics:phyx-mini-0819:phyxminiproblems-problemphyxmini0819-lengthinmeters}
  \lean{PhyXMiniProblems.ProblemPhyXMini0819.lengthInMeters}
  \uses{def:physics:phyx-mini-0819:phyxminiproblems-problemphyxmini0819-lengthquantity}
  Read a physical length in coherent-SI metres
\end{definition}
\begin{proof}
  This is the declared model definition of length In Meters.
\end{proof}
\begin{definition}[angular Speed In Radians Per Second]
  \label{def:physics:phyx-mini-0819:phyxminiproblems-problemphyxmini0819-angularspeedinradianspersecond}
  \lean{PhyXMiniProblems.ProblemPhyXMini0819.angularSpeedInRadiansPerSecond}
  \uses{def:physics:phyx-mini-0819:phyxminiproblems-problemphyxmini0819-angularspeedquantity}
  Read angular speed in radians per coherent-SI second
\end{definition}
\begin{proof}
  This is the declared model definition of angular Speed In Radians Per Second.
\end{proof}
\begin{definition}[moment Of Inertia In Kilogram Meters Squared]
  \label{def:physics:phyx-mini-0819:phyxminiproblems-problemphyxmini0819-momentofinertiainkilogrammeterssquared}
  \lean{PhyXMiniProblems.ProblemPhyXMini0819.momentOfInertiaInKilogramMetersSquared}
  \uses{def:physics:phyx-mini-0819:phyxminiproblems-problemphyxmini0819-momentofinertiaquantity}
  Read an axial moment of inertia in coherent-SI kg m²
\end{definition}
\begin{proof}
  This is the declared model definition of moment Of Inertia In Kilogram Meters Squared.
\end{proof}
\begin{definition}[energy In Joules]
  \label{def:physics:phyx-mini-0819:phyxminiproblems-problemphyxmini0819-energyinjoules}
  \lean{PhyXMiniProblems.ProblemPhyXMini0819.energyInJoules}
  \uses{def:physics:phyx-mini-0819:phyxminiproblems-problemphyxmini0819-energyquantity}
  Read a physical energy in joules, calibrated by Physlib's joule
\end{definition}
\begin{proof}
  This is the declared model definition of energy In Joules.
\end{proof}
\begin{definition}[Disk Label]
  \label{def:physics:phyx-mini-0819:phyxminiproblems-problemphyxmini0819-disklabel}
  \lean{PhyXMiniProblems.ProblemPhyXMini0819.DiskLabel}
  The three labeled small disks in the primary figure
\end{definition}
\begin{proof}
  This is the declared model definition of Disk Label.
\end{proof}
\begin{definition}[Strut Label]
  \label{def:physics:phyx-mini-0819:phyxminiproblems-problemphyxmini0819-strutlabel}
  \lean{PhyXMiniProblems.ProblemPhyXMini0819.StrutLabel}
  The three lightweight struts forming the triangular frame
\end{definition}
\begin{proof}
  This is the declared model definition of Strut Label.
\end{proof}
\begin{definition}[Axis Label]
  \label{def:physics:phyx-mini-0819:phyxminiproblems-problemphyxmini0819-axislabel}
  \lean{PhyXMiniProblems.ProblemPhyXMini0819.AxisLabel}
  The two axes explicitly labeled in the figure
\end{definition}
\begin{proof}
  This is the declared model definition of Axis Label.
\end{proof}
\begin{definition}[Axis Orientation]
  \label{def:physics:phyx-mini-0819:phyxminiproblems-problemphyxmini0819-axisorientation}
  \lean{PhyXMiniProblems.ProblemPhyXMini0819.AxisOrientation}
  Orientation of an axis relative to the plane of the triangular frame
\end{definition}
\begin{proof}
  This is the declared model definition of Axis Orientation.
\end{proof}
\begin{definition}[Small Disk Model]
  \label{def:physics:phyx-mini-0819:phyxminiproblems-problemphyxmini0819-smalldiskmodel}
  \lean{PhyXMiniProblems.ProblemPhyXMini0819.SmallDiskModel}
  Idealization of each small disk when no disk radius is supplied
\end{definition}
\begin{proof}
  This is the declared model definition of Small Disk Model.
\end{proof}
\begin{definition}[Strut Mass Model]
  \label{def:physics:phyx-mini-0819:phyxminiproblems-problemphyxmini0819-strutmassmodel}
  \lean{PhyXMiniProblems.ProblemPhyXMini0819.StrutMassModel}
  Idealization justified by the description of each strut as lightweight
\end{definition}
\begin{proof}
  This is the declared model definition of Strut Mass Model.
\end{proof}
\begin{definition}[Three Disk Figure]
  \label{def:physics:phyx-mini-0819:phyxminiproblems-problemphyxmini0819-threediskfigure}
  \lean{PhyXMiniProblems.ProblemPhyXMini0819.ThreeDiskFigure}
  \uses{def:physics:phyx-mini-0819:phyxminiproblems-problemphyxmini0819-disklabel, def:physics:phyx-mini-0819:phyxminiproblems-problemphyxmini0819-strutlabel, def:physics:phyx-mini-0819:phyxminiproblems-problemphyxmini0819-axislabel, def:physics:phyx-mini-0819:phyxminiproblems-problemphyxmini0819-axisorientation}
  Literal and qualitative information transcribed from image 819.png
\end{definition}
\begin{proof}
  This is the declared model definition of Three Disk Figure.
\end{proof}
\begin{definition}[Three Disk Machine Part]
  \label{def:physics:phyx-mini-0819:phyxminiproblems-problemphyxmini0819-threediskmachinepart}
  \lean{PhyXMiniProblems.ProblemPhyXMini0819.ThreeDiskMachinePart}
  \uses{def:physics:phyx-mini-0819:phyxminiproblems-problemphyxmini0819-massquantity, def:physics:phyx-mini-0819:phyxminiproblems-problemphyxmini0819-lengthquantity, def:physics:phyx-mini-0819:phyxminiproblems-problemphyxmini0819-angularspeedquantity, def:physics:phyx-mini-0819:phyxminiproblems-problemphyxmini0819-momentofinertiaquantity, def:physics:phyx-mini-0819:phyxminiproblems-problemphyxmini0819-energyquantity, def:physics:phyx-mini-0819:phyxminiproblems-problemphyxmini0819-disklabel, def:physics:phyx-mini-0819:phyxminiproblems-problemphyxmini0819-strutlabel, def:physics:phyx-mini-0819:phyxminiproblems-problemphyxmini0819-axislabel, def:physics:phyx-mini-0819:phyxminiproblems-problemphyxmini0819-smalldiskmodel, def:physics:phyx-mini-0819:phyxminiproblems-problemphyxmini0819-strutmassmodel, def:physics:phyx-mini-0819:phyxminiproblems-problemphyxmini0819-threediskfigure}
  Independent physical quantities for the machine part. In particular, rotationalKineticEnergy is not defined from any answer choice. The RigidBody 3 fields expose Physlib's SI-coordinate mechanics, while the corresponding physical scalar magnitudes remain dimensionful
\end{definition}
\begin{proof}
  This is the declared model definition of Three Disk Machine Part.
\end{proof}
\begin{definition}[Matches Primary Figure]
  \label{def:physics:phyx-mini-0819:phyxminiproblems-problemphyxmini0819-matchesprimaryfigure}
  \lean{PhyXMiniProblems.ProblemPhyXMini0819.MatchesPrimaryFigure}
  \uses{def:physics:phyx-mini-0819:phyxminiproblems-problemphyxmini0819-massinkilograms, def:physics:phyx-mini-0819:phyxminiproblems-problemphyxmini0819-lengthinmeters, def:physics:phyx-mini-0819:phyxminiproblems-problemphyxmini0819-disklabel, def:physics:phyx-mini-0819:phyxminiproblems-problemphyxmini0819-strutlabel, def:physics:phyx-mini-0819:phyxminiproblems-problemphyxmini0819-axislabel, def:physics:phyx-mini-0819:phyxminiproblems-problemphyxmini0819-threediskmachinepart}
  All metric and qualitative information read from the primary bitmap
\end{definition}
\begin{proof}
  This is the declared model definition of Matches Primary Figure.
\end{proof}
\begin{definition}[Matches Problem Data]
  \label{def:physics:phyx-mini-0819:phyxminiproblems-problemphyxmini0819-matchesproblemdata}
  \lean{PhyXMiniProblems.ProblemPhyXMini0819.MatchesProblemData}
  \uses{def:physics:phyx-mini-0819:phyxminiproblems-problemphyxmini0819-angularspeedinradianspersecond, def:physics:phyx-mini-0819:phyxminiproblems-problemphyxmini0819-disklabel, def:physics:phyx-mini-0819:phyxminiproblems-problemphyxmini0819-strutlabel, def:physics:phyx-mini-0819:phyxminiproblems-problemphyxmini0819-threediskmachinepart}
  Problem-text information for the requested motion. The point-mass model is the standard interpretation of "small disks" when no disk radii are given; the strut model records the explicit "lightweight" assumption
\end{definition}
\begin{proof}
  This is the declared model definition of Matches Problem Data.
\end{proof}
\begin{definition}[Satisfies Three Disk Rotational Laws]
  \label{def:physics:phyx-mini-0819:phyxminiproblems-problemphyxmini0819-satisfiesthreediskrotationallaws}
  \lean{PhyXMiniProblems.ProblemPhyXMini0819.SatisfiesThreeDiskRotationalLaws}
  \uses{def:physics:phyx-mini-0819:phyxminiproblems-problemphyxmini0819-angularspeedquantity, def:physics:phyx-mini-0819:phyxminiproblems-problemphyxmini0819-massinkilograms, def:physics:phyx-mini-0819:phyxminiproblems-problemphyxmini0819-lengthinmeters, def:physics:phyx-mini-0819:phyxminiproblems-problemphyxmini0819-angularspeedinradianspersecond, def:physics:phyx-mini-0819:phyxminiproblems-problemphyxmini0819-momentofinertiainkilogrammeterssquared, def:physics:phyx-mini-0819:phyxminiproblems-problemphyxmini0819-energyinjoules, def:physics:phyx-mini-0819:phyxminiproblems-problemphyxmini0819-threediskmachinepart}
  The point-mass axial inertia sum and its connection to Physlib's rigid-body mass distribution, inertia tensor, and rotational-energy definition. These are general physical laws for arbitrary angular speed and contain neither the 0.456 J derived energy nor the displayed 0.46 J answer
\end{definition}
\begin{proof}
  This is the declared model definition of Satisfies Three Disk Rotational Laws.
\end{proof}
\begin{lemma}[axis1 Moment Of Inertia eq 57 over 1000]
  \label{lem:physics:phyx-mini-0819:phyxminiproblems-problemphyxmini0819-axis1momentofinertia-eq-57-over-1000}
  \lean{PhyXMiniProblems.ProblemPhyXMini0819.axis1MomentOfInertia_eq_57_over_1000}
  \uses{def:physics:phyx-mini-0819:phyxminiproblems-problemphyxmini0819-momentofinertiainkilogrammeterssquared, def:physics:phyx-mini-0819:phyxminiproblems-problemphyxmini0819-threediskmachinepart, def:physics:phyx-mini-0819:phyxminiproblems-problemphyxmini0819-matchesprimaryfigure, def:physics:phyx-mini-0819:phyxminiproblems-problemphyxmini0819-matchesproblemdata, def:physics:phyx-mini-0819:phyxminiproblems-problemphyxmini0819-satisfiesthreediskrotationallaws}
  The three labeled point masses give I₁ = 0.057 kg m²
\end{lemma}
\begin{proof}
  The conclusion follows from the governing relations and figure readouts named in the statement of axis1 Moment Of Inertia eq 57 over 1000.
\end{proof}
\begin{lemma}[rotational Kinetic Energy eq half inertia mul angular Speed sq]
  \label{lem:physics:phyx-mini-0819:phyxminiproblems-problemphyxmini0819-rotationalkineticenergy-eq-half-inertia-mul-angularspeed-sq}
  \lean{PhyXMiniProblems.ProblemPhyXMini0819.rotationalKineticEnergy_eq_half_inertia_mul_angularSpeed_sq}
  \uses{def:physics:phyx-mini-0819:phyxminiproblems-problemphyxmini0819-angularspeedquantity, def:physics:phyx-mini-0819:phyxminiproblems-problemphyxmini0819-angularspeedinradianspersecond, def:physics:phyx-mini-0819:phyxminiproblems-problemphyxmini0819-momentofinertiainkilogrammeterssquared, def:physics:phyx-mini-0819:phyxminiproblems-problemphyxmini0819-energyinjoules, def:physics:phyx-mini-0819:phyxminiproblems-problemphyxmini0819-threediskmachinepart, def:physics:phyx-mini-0819:phyxminiproblems-problemphyxmini0819-satisfiesthreediskrotationallaws}
  For angular velocity along the selected tensor axis, Physlib's tensor contraction reduces to the scalar law K = I ω² / 2
\end{lemma}
\begin{proof}
  The conclusion follows from the governing relations and figure readouts named in the statement of rotational Kinetic Energy eq half inertia mul angular Speed sq.
\end{proof}
\begin{lemma}[axis1 Kinetic Energy At Four eq 57 over 125]
  \label{lem:physics:phyx-mini-0819:phyxminiproblems-problemphyxmini0819-axis1kineticenergyatfour-eq-57-over-125}
  \lean{PhyXMiniProblems.ProblemPhyXMini0819.axis1KineticEnergyAtFour_eq_57_over_125}
  \uses{def:physics:phyx-mini-0819:phyxminiproblems-problemphyxmini0819-energyinjoules, def:physics:phyx-mini-0819:phyxminiproblems-problemphyxmini0819-threediskmachinepart, def:physics:phyx-mini-0819:phyxminiproblems-problemphyxmini0819-matchesprimaryfigure, def:physics:phyx-mini-0819:phyxminiproblems-problemphyxmini0819-matchesproblemdata, def:physics:phyx-mini-0819:phyxminiproblems-problemphyxmini0819-satisfiesthreediskrotationallaws}
  At 4 rad/s, the exact point-mass-model energy is 0.456 J
\end{lemma}
\begin{proof}
  The conclusion follows from the governing relations and figure readouts named in the statement of axis1 Kinetic Energy At Four eq 57 over 125.
\end{proof}
\begin{definition}[Energy Answer Choice]
  \label{def:physics:phyx-mini-0819:phyxminiproblems-problemphyxmini0819-energyanswerchoice}
  \lean{PhyXMiniProblems.ProblemPhyXMini0819.EnergyAnswerChoice}
  Labels of the four energy choices displayed in the exercise
\end{definition}
\begin{proof}
  This is the declared model definition of Energy Answer Choice.
\end{proof}
\begin{definition}[displayed Energy Joules]
  \label{def:physics:phyx-mini-0819:phyxminiproblems-problemphyxmini0819-displayedenergyjoules}
  \lean{PhyXMiniProblems.ProblemPhyXMini0819.displayedEnergyJoules}
  \uses{def:physics:phyx-mini-0819:phyxminiproblems-problemphyxmini0819-energyanswerchoice}
  The numerical joule value printed beside each answer label
\end{definition}
\begin{proof}
  This is the declared model definition of displayed Energy Joules.
\end{proof}
\begin{definition}[recorded Answer Choice]
  \label{def:physics:phyx-mini-0819:phyxminiproblems-problemphyxmini0819-recordedanswerchoice}
  \lean{PhyXMiniProblems.ProblemPhyXMini0819.recordedAnswerChoice}
  \uses{def:physics:phyx-mini-0819:phyxminiproblems-problemphyxmini0819-energyanswerchoice}
  Dataset metadata recording the supplied answer label
\end{definition}
\begin{proof}
  This is the declared model definition of recorded Answer Choice.
\end{proof}
\begin{definition}[Rounds To Nearest Hundredth]
  \label{def:physics:phyx-mini-0819:phyxminiproblems-problemphyxmini0819-roundstonearesthundredth}
  \lean{PhyXMiniProblems.ProblemPhyXMini0819.RoundsToNearestHundredth}
  displayed is the nearest-hundredth rounding bin containing exact
\end{definition}
\begin{proof}
  This is the declared model definition of Rounds To Nearest Hundredth.
\end{proof}
\begin{definition}[Is Closest Energy Choice]
  \label{def:physics:phyx-mini-0819:phyxminiproblems-problemphyxmini0819-isclosestenergychoice}
  \lean{PhyXMiniProblems.ProblemPhyXMini0819.IsClosestEnergyChoice}
  \uses{def:physics:phyx-mini-0819:phyxminiproblems-problemphyxmini0819-energyanswerchoice, def:physics:phyx-mini-0819:phyxminiproblems-problemphyxmini0819-displayedenergyjoules}
  A choice is at least as close to the exact energy as every other choice
\end{definition}
\begin{proof}
  This is the declared model definition of Is Closest Energy Choice.
\end{proof}
\begin{definition}[Is Unique Closest Energy Choice]
  \label{def:physics:phyx-mini-0819:phyxminiproblems-problemphyxmini0819-isuniqueclosestenergychoice}
  \lean{PhyXMiniProblems.ProblemPhyXMini0819.IsUniqueClosestEnergyChoice}
  \uses{def:physics:phyx-mini-0819:phyxminiproblems-problemphyxmini0819-energyanswerchoice, def:physics:phyx-mini-0819:phyxminiproblems-problemphyxmini0819-isclosestenergychoice}
  A displayed answer is the unique closest energy choice
\end{definition}
\begin{proof}
  This is the declared model definition of Is Unique Closest Energy Choice.
\end{proof}
% --- Archon declaration topology end ---
% --- Archon physics formalization source end ---
